% GENERATED FILE. Edit canonical lessons, then run npm run generate:books.
% canonical-source-hash: fnv1a64:9f2cba229bdf543a
% canonical-lessons: ZH-C292-guangbo, ZH-C292-guanzhong, ZH-C292-zhibo, ZH-C292-tingzhong, ZH-C292-fabiao

\chapter{On the Radio}
\label{ch:zh-on-the-radio}
\hlchaptermodality{292}

\begin{chapteropening}
\textbf{By the end of this chapter:} \emph{I can talk about broadcasts and live shows, viewers and listeners, and what someone published.}

Complete the last lesson of chapter 292: I can talk about broadcasts and live shows, viewers and listeners, and what someone published.
\end{chapteropening}

\section[guǎngbō]{guǎngbō — to broadcast; a broadcast}

\label{lesson:ZH-C292-guangbo}



\begin{quote}
Before the new one: say the Mandarin for to break out suddenly, then the Mandarin for to announce officially.
\end{quote}

\subsection*{You'll want to know: guǎngbō}
\emph{guǎngbō} — \textquotedblleft{}to broadcast; a broadcast\textquotedblright{}.

\begin{quote}
\emph{diàntái měitiān zǎoshang guǎngbō xīnwén} — The station broadcasts the news every morning.
\end{quote}

Its characters come later; for now it is a word for the ear.

\begin{cousinweb}
\emph{guǎngbō} is built from \emph{guǎng} 'wide' and \emph{bō} 'to scatter', both inherited from Middle Chinese, \emph{kwangX} and \emph{paH} (Guangyun, in Baxter's transcription), which glosses the first 'big' and 'broad' and the second 'to spread' and 'to sow': to spread far and wide. The level tone of bō is irregular for \emph{paH}.
\end{cousinweb}

\begin{grammarlens}[title={What you've built}]
One more word for telling what happened.
\end{grammarlens}

\subsection*{Your turn}
\begin{itemize}
\raggedright
  \item \emph{Say it:} \emph{guǎngbō}
  \item \emph{Say it:} \emph{guǎngbō}, once more
  \item \emph{Say it:} \emph{diàntái měitiān zǎoshang guǎngbō xīnwén} — The station broadcasts the news every morning.
  \item \emph{Recall:} say the Mandarin for to break out suddenly, then the Mandarin for to announce officially, then say \emph{guǎngbō} again
\end{itemize}

\begin{tcolorbox}[breakable,colback=teal!4,colframe=teal!35!black,title={Before you move on}]
What does guǎngbō mean? (\textquotedblleft{}To broadcast; a broadcast\textquotedblright{}.) Say it once more.
\end{tcolorbox}
\section[guānzhòng]{guānzhòng — viewers; spectators; the people watching a show}

\label{lesson:ZH-C292-guanzhong}



\begin{quote}
Before the new one: say the Mandarin for to announce officially, then the Mandarin for to broadcast on the radio.
\end{quote}

\subsection*{You'll want to know: guānzhòng}
\emph{guānzhòng} — \textquotedblleft{}viewers; spectators; the people watching a show\textquotedblright{}.

\begin{quote}
\emph{guānzhòng dōu hěn xǐhuan zhège jiémù} — The viewers all love this show.
\end{quote}

Its characters come later; for now it is a word for the ear.

\begin{cousinweb}
\emph{guānzhòng} is built from \emph{guān} 'to watch' and \emph{zhòng} 'a crowd', both inherited from Middle Chinese, \emph{kwan} and \emph{tsyuwngH} (Guangyun, in Baxter's transcription), which glosses the first 'to view' and the second 'many', adding that three people make a crowd.
\end{cousinweb}

\begin{grammarlens}[title={What you've built}]
One more word for telling what happened.
\end{grammarlens}

\subsection*{Your turn}
\begin{itemize}
\raggedright
  \item \emph{Say it:} \emph{guānzhòng}
  \item \emph{Say it:} \emph{guānzhòng}, once more
  \item \emph{Say it:} \emph{guānzhòng dōu hěn xǐhuan zhège jiémù} — The viewers all love this show.
  \item \emph{Recall:} say the Mandarin for to announce officially, then the Mandarin for to broadcast on the radio, then say \emph{guānzhòng} again
\end{itemize}

\begin{tcolorbox}[breakable,colback=teal!4,colframe=teal!35!black,title={Before you move on}]
What does guānzhòng mean? (\textquotedblleft{}Viewers; spectators; the people watching a show\textquotedblright{}.) Say it once more.
\end{tcolorbox}
\section[zhíbō]{zhíbō — to broadcast live; a live show}

\label{lesson:ZH-C292-zhibo}



\begin{quote}
Before the new one: say the Mandarin for to broadcast on the radio, then the Mandarin for the people watching a show.
\end{quote}

\subsection*{You'll want to know: zhíbō}
\emph{zhíbō} — \textquotedblleft{}to broadcast live; a live show\textquotedblright{}.

\begin{quote}
\emph{diànshì zhèngzài zhíbō zúqiú bǐsài} — The football match is on live on TV right now.
\end{quote}

Its characters come later; for now it is a word for the ear.

\begin{cousinweb}
\emph{zhíbō} is built from \emph{zhí} 'straight' and \emph{bō} 'to scatter', both inherited from Middle Chinese, \emph{drik} and \emph{paH} (Guangyun, in Baxter's transcription), which glosses the first 'upright' and 'straight'. The level tone of bō is irregular for \emph{paH}.
\end{cousinweb}

\begin{grammarlens}[title={What you've built}]
One more word for telling what happened.
\end{grammarlens}

\subsection*{Your turn}
\begin{itemize}
\raggedright
  \item \emph{Say it:} \emph{zhíbō}
  \item \emph{Say it:} \emph{zhíbō}, once more
  \item \emph{Say it:} \emph{diànshì zhèngzài zhíbō zúqiú bǐsài} — The football match is on live on TV right now.
  \item \emph{Recall:} say the Mandarin for to broadcast on the radio, then the Mandarin for the people watching a show, then say \emph{zhíbō} again
\end{itemize}

\begin{tcolorbox}[breakable,colback=teal!4,colframe=teal!35!black,title={Before you move on}]
What does zhíbō mean? (\textquotedblleft{}To broadcast live; a live show\textquotedblright{}.) Say it once more.
\end{tcolorbox}
\section[tīngzhòng]{tīngzhòng — listeners; the people listening to the radio}

\label{lesson:ZH-C292-tingzhong}



\begin{quote}
Before the new one: say the Mandarin for the people watching a show, then the Mandarin for to broadcast live.
\end{quote}

\subsection*{You'll want to know: tīngzhòng}
\emph{tīngzhòng} — \textquotedblleft{}listeners; the people listening to the radio\textquotedblright{}.

\begin{quote}
\emph{rúguǒ yǒu tīngzhòng dǎ diànhuà, zhǔchírén jiù huì huídá} — If a listener calls in, the host will answer.
\end{quote}

Its characters come later; for now it is a word for the ear.

\begin{cousinweb}
\emph{tīngzhòng} is built from \emph{tīng} 'to listen' and \emph{zhòng} 'a crowd', both inherited from Middle Chinese, \emph{theng} and \emph{tsyuwngH} (Guangyun, in Baxter's transcription), which glosses the first 'to listen' and the second 'many'.
\end{cousinweb}

\begin{grammarlens}[title={What you've built}]
One more word for telling what happened.
\end{grammarlens}

\subsection*{Your turn}
\begin{itemize}
\raggedright
  \item \emph{Say it:} \emph{tīngzhòng}
  \item \emph{Say it:} \emph{tīngzhòng}, once more
  \item \emph{Say it:} \emph{rúguǒ yǒu tīngzhòng dǎ diànhuà, zhǔchírén jiù huì huídá} — If a listener calls in, the host will answer.
  \item \emph{Recall:} say the Mandarin for the people watching a show, then the Mandarin for to broadcast live, then say \emph{tīngzhòng} again
\end{itemize}

\begin{tcolorbox}[breakable,colback=teal!4,colframe=teal!35!black,title={Before you move on}]
What does tīngzhòng mean? (\textquotedblleft{}Listeners; the people listening to the radio\textquotedblright{}.) Say it once more.
\end{tcolorbox}
\section[fābiǎo]{fābiǎo — to publish an article; to make public}

\label{lesson:ZH-C292-fabiao}



\begin{quote}
Before the new one: say the Mandarin for to broadcast live, then the Mandarin for the people listening to the radio.
\end{quote}

\subsection*{You'll want to know: fābiǎo}
\emph{fābiǎo} — \textquotedblleft{}to publish an article; to make public\textquotedblright{}.

\begin{quote}
\emph{tā fābiǎo le yì piān xiǎoshuō} — He published a short story.
\end{quote}

Its characters come later; for now it is a word for the ear.

\begin{cousinweb}
\emph{fābiǎo} is built from \zh{发} 'to send out' and \emph{biǎo} 'the outside', both inherited from Middle Chinese, \emph{pjot} and \emph{pjewX} (Guangyun, in Baxter's transcription), which glosses the second 'to make clear' and 'a letter to the throne': to send out where all can see.
\end{cousinweb}

\begin{grammarlens}[title={What you've built}]
One more word for telling what happened.
\end{grammarlens}

\subsection*{Your turn}
\begin{itemize}
\raggedright
  \item \emph{Say it:} \emph{fābiǎo}
  \item \emph{Say it:} \emph{fābiǎo}, once more
  \item \emph{Say it:} \emph{tā fābiǎo le yì piān xiǎoshuō} — He published a short story.
  \item \emph{Recall:} say the Mandarin for to broadcast live, then the Mandarin for the people listening to the radio, then say \emph{fābiǎo} again
\end{itemize}

\begin{tcolorbox}[breakable,colback=teal!4,colframe=teal!35!black,title={Before you move on}]
What does fābiǎo mean? (\textquotedblleft{}To publish an article; to make public\textquotedblright{}.) Say it once more.
\end{tcolorbox}
